\documentclass{article}
\usepackage[utf8]{inputenc}
\usepackage[spanish]{babel}
\usepackage{amsmath,amssymb}
\begin{document}
\section{Introducción}
La función de distribución, según el teorema, cumple que $f(x) = x^2 + 2x - 1$ para todo $x \in \mathbb{R}$. ¿Qué ocurre si $\alpha \leq \beta$? ¡Es importante!
Además, la integral $\int_0^1 e^{-x^2}\,dx \approx 0{,}7468$ y la serie converge.
\begin{equation}
\sum_{i=1}^{n} a_i = \frac{n(n+1)}{2}
\end{equation}
\begin{equation}
E = mc^2 \label{eq:e}
\end{equation}
Sea $\varepsilon > 0$; entonces existe $\delta > 0$ tal que $|x - a| < \delta \Rightarrow |f(x) - L| < \varepsilon$.
\begin{itemize}
\item Primer punto con $\sqrt{2}$ y $\infty$.
\item Segundo punto: $A \subseteq B \cup C$.
\end{itemize}
Una palabra muy larga para forzar la separación silábica: electroencefalografistas anticonstitucionalmente desoxirribonucleico.\footnote{Nota al pie.}
\[ \lim_{x \to 0} \frac{\sin x}{x} = 1 \]
\end{document}
